\documentclass[../../../../main.tex]{subfiles}
\begin{document}

3辺の長さが $BC=2a$，$CA=2b$，$AB=2c$ であるような鋭角三角形 $\triangle ABC$ の
3辺 $BC$，$CA$，$AB$ の中点をそれぞれ $L$，$M$，$N$ とする．
線分 $LM$，$MN$，$NL$ に沿って三角形を折り曲げ，四面体をつくる．
その際，線分 $BL$ と $CL$，$CM$ と $AM$，$AN$ と $BN$ はそれぞれ同一視されて，
長さが $a$，$b$，$c$ の辺になるものとする．

\begin{figure}[htbp]
  \centering
  \begin{tikzpicture}[scale=1.5]
    \coordinate (A) at (1.2, 2.2);
    \coordinate (B) at (0, 0);
    \coordinate (C) at (3.2, 0);

    \coordinate (L) at ($(B)!0.5!(C)$);
    \coordinate (M) at ($(C)!0.5!(A)$);
    \coordinate (N) at ($(A)!0.5!(B)$);

    \coordinate (P) at ($(M)!0.5!(N)$);
    \coordinate (Q) at ($(B)!0.5!(L)$);

    \draw[thick] (A) -- (B) -- (C) -- cycle;
    \draw[thick] (L) -- (M) -- (N) -- cycle;

    \fill (L) circle (1.5pt) node[below] {$L$};
    \fill (M) circle (1.5pt) node[above right] {$M$};
    \fill (N) circle (1.5pt) node[above left] {$N$};
    \fill (P) circle (1.5pt) node[above] {$P$};
    \fill (Q) circle (1.5pt) node[below] {$Q$};

    \node[above] at (A) {$A$};
    \node[below left] at (B) {$B$};
    \node[below right] at (C) {$C$};

    \node[above left] at ($(A)!0.5!(N)$) {$c$};
    \node[below left] at ($(N)!0.5!(B)$) {$c$};
    \node[above right] at ($(A)!0.5!(M)$) {$b$};
    \node[below right] at ($(M)!0.5!(C)$) {$b$};
    \node[below] at ($(B)!0.5!(L)$) {$a$};
    \node[below] at ($(L)!0.5!(C)$) {$a$};
  \end{tikzpicture}
\end{figure}

\begin{enumerate}
  \item 線分 $MN$，$BL$ の中点をそれぞれ $P$，$Q$ とする．
  四面体を組み立てたとき，空間内の線分 $PQ$ の長さを求めよ．

  \item この四面体の体積を $a$，$b$，$c$ を用いて表せ．
\end{enumerate}

\end{document}
